\section{Part II.a Appendix: Witnesses for Unwitnessed Structures}

\subsection{Witness}

\subsubsection*{Epistemic status}

An external audit found that six structures carrying load-bearing theorems were never
\emph{constructed} anywhere in the development.  A theorem stating ``given a structure
$S$, $\ldots$'' is true but unwitnessed if no object builds an $S$; its physical reading
is then a conditional, not a result.  The most consequential case was $\text{\textsc{ScaleShift}}$,
the carrier of the identification ``$\hbar$ is an exact scale step.''

This appendix builds them.  What that does and does not settle is stated first, because
the distinction is the whole point of the exercise.

\paragraph{What a witness settles:}
\begin{enumerate}
\item The theorems are \emph{not vacuous} — there exists something they apply to;
\item the structure is \emph{compatible with A1} — the same object can carry both, so
  the framework does not forbid it;
\item in the $\text{\textsc{ScaleShift}}$ case, the shift is \emph{exact and finite}, not a
  first-order truncation, which is what the original claim needed.
\end{enumerate}

\paragraph{What a witness does \emph{not} settle:}

It is not a derivation.  Exhibiting one model that satisfies A1 and carries
a scale shift does not show the axioms \emph{force} a scale shift; a different model
of A1 may have none.  So $\hbar$'s identification with the step remains an
identification.  The gap between ``compatible'' and ``forced'' is real and is not
closed here.

\paragraph{The model:}

$\mathbb{R}[X]$ with $d/dX$ is the log-scale axis with its derivation: it satisfies A1
for $n = 1$.  The shift $X \mapsto X + 1$ is a ring homomorphism — this is why the
construction works and why ``multiply by the scale factor'' does not: for a unit
$u$, $(u \cdot x)(u \cdot y) = u^2xy \neq u(xy)$, so multiplication by a scale factor is
additive but not multiplicative and is not a $\text{\textsc{ScaleShift}}$.  Translation along
the scale axis is.

The witness `shift\_is\_exact` shows the step is not infinitesimal: $\Delta X = 1$, a
finite difference, not a derivative.  The witness `shift\_commutes\_with\_derivation` shows
the two structures are compatible rather than merely coexisting.

\subsubsection*{I. The scale axis satisfies A1}

\begin{definition}
\label{def:polyScaleAlgebra}
$\mathbb{R}[X]$ with $d/dX$ is a model of A1 for one direction: the log-scale
axis carrying its own derivation. \lean{Witness.polyScaleAlgebra}

The instance $\text{\textsc{ScaleAlgebra}} \, 1 \, (\mathbb{R}[X])$ is defined by setting the derivation
operator $d$ to be the formal derivative $d/dX$, with the additive and multiplicative
properties inherited from $\mathbb{R}[X]$.
\end{definition}

\subsubsection*{II. And it carries an exact scale shift}

\begin{definition}
\label{def:polyShift}
The witness for $\text{\textsc{ScaleShift}}$: translation by one step along the scale axis.
\lean{Witness.polyShift}

A ring homomorphism, hence a legitimate $\text{\textsc{ScaleShift}}$.  Note what is \emph{not} a
$\text{\textsc{ScaleShift}}$: multiplication by a scale factor $u$, since
$(u \cdot x)(u \cdot y) = u^2xy \neq u(xy)$ — additive but not multiplicative.  The shift
is the translation, not the rescaling.

Formally, the shift operator $T$ acts on $p \in \mathbb{R}[X]$ by composition:
$T(p) = p \circ (X + 1)$, and it preserves both addition and multiplication of
polynomials.
\end{definition}

\begin{theorem}
\label{thm:shift_is_exact}
The step is finite, not infinitesimal. \lean{Witness.shift\_is\_exact}

$\Delta X = (X+1) - X = 1$.  So the non-commutativity of $\text{Crossed.commutator\_eq}$
produces is exact, which is what the original claim required and what a
first-order deformation could not give.

\begin{proof}
The finite difference $\Delta X = T(X) - X = (X + 1) - X = 1$ is computed by
direct evaluation: $(X+1) - X = 1$. This is the exact step, not infinitesimal.
\end{proof}
\end{theorem}

\begin{theorem}
\label{thm:shift_nontrivial}
The shift is not the identity, so the commutator it generates is
genuinely nonzero: the witness is not degenerate. \lean{Witness.shift\_nontrivial}

That is, $T(X) \neq X$ where $T(X) = X + 1$.

\begin{proof}
Suppose $T(X) = X$.  Then by evaluating both polynomials at $0$, we have
$\text{eval}(0, T(X)) = \text{eval}(0, X)$, which gives
$\text{eval}(0, X + 1) = \text{eval}(0, X)$, hence $0 + 1 = 0$, a contradiction.
\end{proof}
\end{theorem}

\begin{theorem}
\label{thm:shift_commutes_with_derivation}
The two structures are compatible, not merely coexisting: translation
commutes with differentiation, so the shift moves along the scale axis without
disturbing A1's derivation. \lean{Witness.shift\_commutes\_with\_derivation}

For all $p \in \mathbb{R}[X]$,
\[
\frac{d}{dX} T(p) = T\left(\frac{d}{dX} p\right).
\]

\begin{proof}
The derivative commutes with composition: $\frac{d}{dX}(p \circ (X+1)) =
\left(\frac{dp}{dX}\right) \circ (X+1)$ by the chain rule. This shows
$\frac{d}{dX}(T(p)) = T\left(\frac{d}{dX}(p)\right)$ for all polynomials.
\end{proof}
\end{theorem}

\subsubsection*{III. The remaining five structures}

Each was flagged as never constructed.  Witnesses are cheap for four of them
and are given so the theorems are known non-vacuous; the fifth is discussed.

\begin{definition}
\label{def:matrixTrace}
$\text{\textsc{Trace}}$: the matrix trace is one, and it is the canonical example — the
`cyclic' field is exactly $\text{trace}(AB) = \text{trace}(BA)$.
Valued in the ring itself via the scalar embedding. \lean{Witness.matrixTrace}

For $n \in \mathbb{N}$, the trace witness on $n \times n$ matrices over $\mathbb{R}$ is
defined by $\tau(M) = \text{trace}(M) \cdot 1$ where $1$ is the identity matrix in
$\text{Matrix}(\text{Fin} \, n, \text{Fin} \, n, \mathbb{R})$. The cyclic property follows
from the cyclic property of the matrix trace.
\end{definition}

\begin{definition}
\label{def:adTransport}
$\text{\textsc{DirTransport}}$: the adjoint action of a ring on itself.  Transport along
$x$ is $[\cdot, x]$ (the commutator), which is additive in the direction, additive in the value, and
a derivation — the three fields, in order. \lean{Witness.adTransport}

For a ring $M$, the adjoint transport is defined by $D(x, m) = [x, m] = xm - mx$,
the Lie bracket (commutator). This is additive in $x$, additive in $m$, and satisfies
the Leibniz rule: $[x, mm'] = [x,m]m' + m[x,m']$.
\end{definition}

\begin{definition}
\label{def:trivialFlow}
$\text{\textsc{ScaleFlow}}$: the trivial flow, which fixes everything.  A witness that
the structure is inhabited; it is also the degenerate case in which every
coupling is a fixed point, so it shows the definitions are consistent without
pretending to model running. \lean{Witness.trivialFlow}

For any type $S$, the trivial flow is defined by $\text{flow}(t, g) = g$ for all $t$
and $g \in S$. This trivially satisfies $\text{flow}(0, g) = g$ and
$\text{flow}(t + u, g) = \text{flow}(t, \text{flow}(u, g))$.
\end{definition}

\begin{definition}
\label{def:linearFlow}
A non-degenerate witness as well: translation of a real coupling by the
log-scale, which is the linear running of $\text{\texttt{Running.lean}}$ in flow form.
\lean{Witness.linearFlow}

The linear flow on $\mathbb{R}$ is defined by $\text{flow}(t, g) = g + t$ for $t \in \mathbb{R}$
and $g \in \mathbb{R}$. This is the additive translation along the real line scaled by the
log-scale parameter $t$.
\end{definition}

\begin{theorem}
\label{thm:linearFlow_moves}
The linear flow has no fixed point at all, so the fixed-point theorems
of $\text{\texttt{RG.lean}}$ are not about an empty case: a non-degenerate witness.
\lean{Witness.linearFlow\_moves}

For all $g \in \mathbb{R}$, $\text{flow}(1, g) \neq g$.

\begin{proof}
By the definition of the linear flow, $\text{flow}(1, g) = g + 1 \neq g$ for any $g \in \mathbb{R}$,
which follows from $1 \neq 0$ in $\mathbb{R}$.
\end{proof}
\end{theorem}

\begin{theorem}
\label{thm:constMoment_conserved}
$\text{\textsc{Conserved}}$ (from $\text{\texttt{Waves.lean}}$): a constant moment.  This is the witness
that the monopole and dipole theorems are non-vacuous — and it is also exactly
the physical content, since conservation \emph{means} the moment is constant.
\lean{Witness.constMoment\_conserved}

For any constant $c \in \mathbb{R}$, the function $\mu(\cdot) = c$ is $\text{\textsc{Conserved}}$.

\begin{proof}
A constant function has derivative zero, so $\frac{d}{dt}(c) = 0 = 0 \cdot 1$ for all time.
Thus the hasDerivAt condition is satisfied for any constant by the constant derivative theorem.
\end{proof}
\end{theorem}

\begin{definition}
\label{def:constantDrift}
$\text{\textsc{CosmicHistory}}$: a history whose drift is a real multiple of $1$, which
is spatially constant because derivations kill constants.  The `drift' field is
therefore satisfiable rather than an unfulfillable demand. \lean{Witness.constantDrift}

Given a base polynomial $\text{base} \in \mathbb{R}[X]$ and a direction $n = 1$, the constant-drift
$\text{\textsc{CosmicHistory}}$ is defined by $\text{drift}(t) = C(t)$, the constant polynomial
in $X$ with coefficient $t$. This satisfies spatial constancy because
$\frac{d}{dX}(C(t)) = 0$ for any constant polynomial.
\end{definition}

\subsubsection*{IV. What is now settled, and what is not}

\begin{theorem}
\label{thm:scaleShift_is_witnessed}
Collected: the flagship structure is inhabited by an object satisfying
A1, and the step it produces is exact. \lean{Witness.scaleShift\_is\_witnessed}

$\Delta X = 1$ and $T(X) \neq X$ where $T$ is the scale shift witness.

That removes ``unwitnessed'' from $\text{\textsc{ScaleShift}}$.  It does \emph{not} remove
``identification'': one model carrying both structures does not show every model
of A1 must carry a scale shift, so $\hbar = (\text{the step})$ remains an identification and
stays registered as such.

\begin{proof}
This combines the two key properties: $\text{shift\_is\_exact}$ establishes
$\Delta X = 1$, and $\text{shift\_nontrivial}$ establishes $T(X) \neq X$.
Both follow from the definition of $T(p) = p \circ (X+1)$ applied to the generator $X$.
\end{proof}
\end{theorem}
